\documentclass[12pt]{article}
\usepackage{auck2e,multicol,amssymb,amsmath,color}
\usepackage{epsfig}
\usepackage[english]{babel}
\usepackage{blindtext}
\usepackage{microtype}
\usepackage{minted}

%\usepackage{graphicx}
\newcommand{\ds}{\displaystyle}
\newcommand{\mmark}[1]{[{#1 mark(s)}]}
\newcommand{\markk}[1]{[{\bf #1 marks}]}
\newcommand{\dspace}{\renewcommand{\baselinestretch}{2} \small \normalsize}
\newcommand{\sspace}{\renewcommand{\baselinestretch}{1} \small \normalsize}
\newcommand{\btheta}{\mbox{\boldmath $\theta$}}
\newcommand{\blambda}{\mbox{\boldmath $\lambda$}}
\newcommand{\bTheta}{\mbox{\boldmath $\Theta$}}
\newcommand{\bzeta}{\mbox{\boldmath $\zeta$}}
\newcommand{\bpi}{\mbox{\boldmath $\pi$}}
\newcommand{\bphi}{\mbox{\boldmath $\phi$}}
\newcommand{\ba}{\mbox{\boldmath $a$}}
\newcommand{\bb}{\mbox{\boldmath $b$}}
\newcommand{\bd}{\mbox{\boldmath $d$}}
\newcommand{\bm}{\mbox{\boldmath $m$}}
\newcommand{\bx}{\mbox{\boldmath $x$}}
\newcommand{\bh}{\mbox{\boldmath $h$}}
\newcommand{\bu}{\mbox{\boldmath $u$}}
\newcommand{\bP}{\mbox{\boldmath $P$}}
\newcommand{\bQ}{\mbox{\boldmath $Q$}}
\newcommand{\bR}{\mbox{\boldmath $R$}}
\newcommand{\by}{\mbox{\boldmath $y$}}
\newcommand{\bY}{\mbox{\boldmath $Y$}}
\newcommand{\bX}{\mbox{\boldmath $X$}}
\newcommand{\bU}{\mbox{\boldmath $U$}}
\newcommand{\bD}{\mbox{\boldmath $D$}}
\newcommand{\bV}{\mbox{\boldmath $V$}}
\newcommand{\bH}{\mbox{\boldmath $H$}}
\newcommand{\bgamma}{\mbox{\boldmath $\gamma$}}
\newcommand{\bGamma}{\mbox{\boldmath $\Gamma$}}
\newcommand{\balpha}{\mbox{\boldmath $\alpha$}}
\newcommand{\bbeta}{\mbox{\boldmath $\beta$}}
\newcommand{\bdelta}{\mbox{\boldmath $\delta$}}
\newcommand{\bomega}{\mbox{\boldmath $\omega$}}
\newcommand{\bpsi}{\mbox{\boldmath $\psi$}}
\newcommand{\bepsilon}{\mbox{\boldmath $\epsilon$}}
\newcommand{\be}{\begin{eqnarray}}
\newcommand{\ee}{\end{eqnarray}}
\def\BackgroundEPS#1#2#3#4{%
    \special{ps: @beginspecial @setspecial initmatrix
                0.9 setgray #2 #3 translate #4 dup scale}
    \special{ps: plotfile #1}
    \special{ps: @endspecial}
}
\newcommand{\wbugs}{{\sffamily\textbf{WinBUGS }}}
\newcommand{\coda}{{\sffamily\textbf{CODA }}}
\newcommand{\glm}{{\sffamily\textbf{G.L.M }}}
\newcommand{\rwb}{{\sffamily\textbf{R2WinBUGS }}}
\newcommand{\rar}{{\sffamily\textbf{Rock and roll man! }}}
\newcommand{\anova}{{\sffamily\textbf{ANOVA }}}
\newcommand{\mcmc}{{\sffamily\textbf{MCMC }}}
\newcommand{\dic}{{\sffamily\textbf{DIC }}}
\newcommand{\dagr}{{\sffamily\textbf{DAG }}}
\newcommand{\boxeq}[1]{%
\fboxrule=0.5mm%
\fboxsep=5mm%
\begin{center}\fcolorbox{black}{yellow}{\begin{minipage}{0.9\linewidth}
\begin{Beqnarray} #1 \end{Beqnarray}\end{minipage}}\end{center}%
\fboxrule=1pt%
\fboxsep=3pt%
}

\newcommand{\boxst}[1]{%
\fboxrule=0.5mm%
\fboxsep=5mm%
\begin{center}\fcolorbox{black}{yellow}{ \begin{minipage}{0.9\linewidth}#1 \end{minipage} }\end{center}%
\fboxrule=1pt%
\fboxsep=3pt%
}
%%%%%%%%%%%%%%%%%%%%%%%%%%%%%%%%%%%%%%%%%%%%%%%%%%%%%%%%%%%%%%%%%%%%%%%%
% Macro         :   \BackgroundEPS{filename}{tx}{ty}{s}
%-----------------------------------------------------------------------
% Description   :   Display an EPS file in the background of the
%                   current page.
%-----------------------------------------------------------------------
% Author        :   Bernard Desruisseaux (bernard@XYZZYcs.mcgill.ca)
%-----------------------------------------------------------------------
% Usage         :   \BackgroundEPS{pic.eps}{80}{100}{2}
%                   \BackgroundEPS{"`gunzip -c pic.eps.gz"}{80}{100}{2}
%                   \BackgroundEPS{"`zcat pic.eps.Z"}{80}{100}{2}
%%%%%%%%%%%%%%%%%%%%%%%%%%%%%%%%%%%%%%%%%%%%%%%%%%%%%%%%%%%%%%%%%%%%%%%%

\newcommand{\given}{\,|\,}

\paper{Introduction to Bayesian Statistics}
\due{Not for assessment}
\thetitle{\textbf{Tutorial 1}} 
\begin{document}


%\BackgroundEPS{bayes.eps}{0}{0}{1.2}
\makeband
You may do the questions with a pen and paper (you may need a calculator),
or in R or any other language of your choice.

\begin{Questions}

\item 
A customs officer inspects a package containing five dolls, and thinks
that the dolls might contain hidden drugs. Let $\theta$ be the number of dolls
that contain illegal drugs, so the possible values are 0, 1, 2, 3, 4, or 5.
The customs officer has prior probabilities
$0.5, 0.025, 0.025, 0.025, 0.025$, and $0.4$ for these six possibilities
respectively.

\begin{enumerate}
\item [(a)] The prior probabilities I wrote down are not entirely arbitrary.
They are also not a simple uniform distribution, but use uniform
assumptions in a couple of ways. See if you can identify them.

\item [(b)] If using R or another programming language,
create {\tt theta} and {\tt prior} vectors to represent the first two
columns of a Bayes Box. Otherwise, just write down the columns on paper.
\end{enumerate}

The customs officer does not want to damage all of the dolls by inspecting them
all. Instead, he chooses one
arbitrarily and thoroughly inspects it, finding no drugs. Let $D$
be the data that the inspected doll contained no drugs.

\begin{enumerate}
\item[(c)] What are the values of the likelihoods?
\item[(d)] Find the posterior distribution for {\tt theta} and hence the posterior
probability that there are drugs.
\end{enumerate}


\item 
Let $x_i$ be the number of times a new puppy pees on the floor on
day $i$. Given a rate parameter $\lambda$, this has a Poisson
distribution with mean $\lambda$:
\begin{align}
p(x_i \given \lambda) &= \frac{\lambda^{x_i}\exp(-\lambda)}{x_i!}.
\end{align}
Four days are observed, and the data is
$\boldsymbol{x} = \{x_1, x_2, x_3, x_4\} = \{2, 0, 2, 3\}$.
This one is best done in a computer.

\begin{enumerate}
\item [(a)] Create a vector for a uniform {\tt prior}
for {\tt lambda} between 0 and 10, in steps of 0.1.
Note: We are approximating a continuous parameter using a fine grid.
\item [(b)] Calculate the likelihoods. You can use the product
of the Poisson formula four times, or (in R at least) you can
use \mintinline{r}{dpois()}.
\item [(c)] Calculate the posterior distribution and plot it.
\item [(d)] Find the expected value of $\lambda$ according to the
posterior distribution.
\item [(e)] The classical maximum likelihood estimate is equal
to the sample mean $\bar{x}$. How different is this from the answer
to (d). Why is it different? There are two main reasons.
\end{enumerate}





\item

A bad joke suggests that a Bayesian is ``one who,
vaguely expecting a horse, and catching a glimpse of a donkey, strongly believes he has seen a mule''. The joke comes from the fact that,
sometimes (but not always),
if the prior and likelihood peak in different places, the
posterior peaks in the middle of the two.

Assume there is a parameter $\theta$ whose prior is a standard normal
distribution, and that some data gave a gaussian likelihood function
with a peak $\theta=20$:
\begin{align}
p(\theta)   &\propto \exp \left(-\frac{1}{2}\theta^2\right) \\
p(x \given \theta) &\propto \exp \left(-\frac{1}{2}(\theta - 20)^2\right)
\end{align}

\begin{enumerate}
\item [(a)] Plot the prior, likelihood, and posterior for $\theta$ on the
                same set of axes. You will see the origin of the bad joke
                ($\theta\approx 0$ represents a horse, $\theta\approx 20$
                 represents a donkey, and
                 $\theta \approx 10$ represents a mule).
\item [(b)] Replace the prior with a Cauchy distribution\footnote{This is also
a $t$ distribution with shape parameter $\nu=1$.}
                $p(\theta) \propto 1/(1+\theta^2)$, which has much heavier tails
                than a normal distribution. Re-make the plot. Does
                the joke still work?

\item [(c)] Is the marginal likelihood higher with the normal
prior or the Cauchy prior? Why?
\end{enumerate}


%\item \textbf{Task:} Solve the ``taxi'' problem from lectures. First you'll
%reproduce my results from lectures with a slight change allowing for the
%possibility of more than 100 taxis (you'll go up to 1000 instead). Then you'll
%calculate what happens with a bigger dataset.

%\begin{enumerate}
%\item Make a vector of possible values for $N$ (the number of taxis in the city)
%ranging from 1 to 1000, and also make a vector to represent a uniform prior.
%\item Calculate the posterior distribution for $N$ given the example data from
%class (the observed taxi number was $x=42$).
%\item Plot the posterior distribution.
%\item To get more information about the number of taxis in the city, you decide
%to gather more data by sleeping out for a further nine nights. Altogether, the
%ten taxi numbers you see are $\{42, 55, 12, 77, 8, 80, 2, 22, 19, 50\}$. If
%you want to infer $N$ from this data, what will the likelihood be? (Hint: this
%will involve the product of 10 terms).
%\item Calculate and plot the posterior distribution for $N$ given the bigger
%data set. Does the result make sense to you?
%\end{enumerate}

\end{Questions}
%\rar


\end{document}




